\section{Introdução}

O objetivo deste trabalho é analisar uma série temporal mensal e ajustar modelos ARIMA ou SARIMA adequados, conforme solicitado no enunciado do Trabalho Aplicado 1 da disciplina Análise de Séries Temporais. A série escolhida foi o índice de volume de vendas da Pesquisa Mensal de Comércio, aqui representado pela variável \texttt{pmc\_volume\_index}.

A análise segue um fluxo de Box--Jenkins: descrição da série, investigação de estacionariedade, ajuste de modelos concorrentes, avaliação por critérios de informação, diagnóstico de resíduos e previsão. Para que a avaliação preditiva seja observável, os últimos 24 meses da amostra foram reservados como conjunto de teste. Assim, as previsões produzidas pelos modelos concorrentes puderam ser comparadas com valores efetivamente observados.

O foco do trabalho é univariado. Embora a base bruta contenha também uma coluna de SELIC, essa variável foi preservada apenas como contexto e não entrou no modelo principal. Essa decisão mantém o trabalho aderente ao objetivo central do enunciado: identificar, ajustar e diagnosticar modelos ARIMA/SARIMA para uma série temporal.

